\section{Coercive equations in presented Hilbert spaces}
\label{sec:hilbert}

This section supplies an infinite-dimensional application of the framework.
Its constructive content is stronger than verifying a proposed solution:
the input data produce a regular solution name with a prescribed
primitive-recursive clock. It also identifies why the corresponding statement
for an arbitrary bounded bilinear form and a weakly presented functional
cannot be asserted without additional information. Computable versions of
Riesz representation and Lax--Milgram have been studied previously
\cite{BrattkaYoshikawa2006}; the argument here is a self-contained quantitative
specialization designed for the arithmetic target.

\subsection{The presentation and its actual inputs}

\begin{definition}[Rational pre-Hilbert presentation]
A rational pre-Hilbert presentation is a coded rational vector space \(H_0\)
with primitive-recursive zero, addition and rational scalar multiplication,
and a rational-valued symmetric positive-definite inner product
\(\langle\ ,\ \rangle\). All vector-space and inner-product laws have their
ordinary finite arithmetic proofs. Its named completion \(H\) consists of
regular names, identified by distance zero. We write
\(\|v\|^2=\langle v,v\rangle\).
\end{definition}

Examples are finite rational sequences with their Euclidean inner product,
or rational polynomials vanishing at \(0,1\), with inner product
\(\int_0^1p'q'\). The dimension of a finite code may grow during an
algorithm. There is no hypothesis that \(H\) is finite-dimensional or totally
bounded.

Norms need not be computed by a real square root. For \(v\in H_0\) and
rational \(r\geq0\), the test \(\|v\|\leq r\) means
\(\langle v,v\rangle\leq r^2\). Rational upper approximations to a norm can
be obtained by a bounded integer square-root calculation. For example,
search integers \(j\) up to
\(\lceil2^t(2+\langle v,v\rangle)\rceil\) for the first
\((j2^{-t})^2\geq\langle v,v\rangle\).
The returned upper radius differs from the norm by at most \(2^{-t}\).

\begin{definition}[Strong operator and vector data]
\label{def:strong-operator}
Fix primitive-recursive maps
\[
 a:H_0\times\N\longrightarrow H_0,\qquad b:\N\longrightarrow H_0.
\]
They name a real-linear operator \(A:H\to H\) and a vector \(b_\infty\in H\)
with errors
\[
 \|a(v,k)-Av\|\leq2^{-k},\qquad
 \|b(k)-b_\infty\|\leq2^{-k}.
 \tag{H1}
\]
Positive rational inputs \(m,M\), with \(m\leq M\), and a rational \(B\geq0\)
certify
\[
 \langle Av,v\rangle\geq m\|v\|^2,\qquad
 \|Av\|\leq M\|v\|,\qquad \|b_\infty\|\leq B.
 \tag{H2}
\]
The maps and their proved certificates, not an unnamed abstract operator, are
the inputs to the arithmetic construction.
\end{definition}

The notation in (H1)--(H2) has a finite reading. One supplies the pairwise
Cauchy estimates for \(a(v,k)\) and \(b(k)\); approximate linearity, for example
\[
\|a(sv+tw,k)-s\,a(v,k)-t\,a(w,k)\|
       \leq(1+|s|+|t|)2^{-k};
\]
and boundedness and coercivity estimates with approximation errors.
For coercivity an adequate entirely rational contract is
\[
\langle a(v,k),v\rangle
 \geq m\langle v,v\rangle-2^{-k}(1+\langle v,v\rangle).
\]
For boundedness one may use every rational upper radius \(r\) of \(v\):
\[
\langle v,v\rangle\leq r^2
 \ \Longrightarrow\ \|a(v,k)\|\leq Mr+2^{-k}.
\]
Here \(r\) is a free numerical parameter. These contracts define the named
operator by extension from the dense space and prove (H2); rational
separation removes auxiliary errors. Conversely, (H1)--(H2) immediately
give these finite contracts. For the exact \(2^{-k}\) approximation in (H1),
one can normalize the supplied Cauchy family by a fixed precision shift.
One must prove the contracts for a concrete presentation; merely storing an
assertion alongside the code is insufficient.

\subsection{The elementary quantitative engine}

\begin{lemma}[Perturbed contraction]
\label{lem:perturb}
Let \(0\leq q<1\), \(\delta=1-q\), and suppose nonnegative quantities satisfy
\[
             e_{j+1}\leq q e_j+\eta,\qquad e_0\leq R .
\]
Then
\[
 e_j\leq q^jR+\eta\sum_{i<j}q^i
       =q^jR+\frac{\eta(1-q^j)}{\delta}
       \leq q^jR+\frac{\eta}{\delta}.
 \tag{H3}
\]
All claims about rational inputs are quantifier-free induction consequences.
\end{lemma}
\begin{proof}
At \(j=0\) the finite sum is zero. Multiplying the inductive estimate by
\(q\geq0\), adding \(\eta\), and reindexing gives the next
estimate. The geometric identity follows by subtracting \(q\) times the
sum from the sum. The denominator is positive. These are finite addition,
multiplication, comparison and induction steps.
\end{proof}

\begin{lemma}[A clock without logarithms]
\label{lem:clock}
If \(0<\delta\leq1\), \(q=1-\delta\), and \(h\in\N\) satisfies \(h\delta\geq1\),
then \(q^h\leq\tfrac12\). Consequently, for integers \(R\geq1\),
\[
       N_n=h(n+R+3)\quad\Longrightarrow\quad
       Rq^{N_n}\leq2^{-n-3}.
 \tag{H4}
\]
\end{lemma}
\begin{proof}
Induction gives
\((1+j\delta)(1-\delta)^j\leq1\).
For the step use
\[
(1+(j+1)\delta)(1-\delta)
=1+j\delta-(j+1)\delta^2\leq1+j\delta.
\]
At \(j=h\), \(1+h\delta\geq2\).
Now \(q^{h(n+R+3)}\leq2^{-(n+R+3)}\).
The elementary inequality \(R\leq2^R\) proves (H4).
\end{proof}

\begin{lemma}[Coercivity supplies the contraction gap]
\label{lem:coercive-contraction}
For the data in Definition~\ref{def:strong-operator}, put
\[
 \lambda=\frac{m}{M^2},\quad
 \delta=\frac{m^2}{2M^2},\quad q=1-\delta,\quad
 T(v)=v-\lambda Av+\lambda b_\infty .
 \tag{H5}
\]
Then \(1/2\leq q<1\) and
\[
                     \|T(v)-T(w)\|\leq q\|v-w\|.
 \tag{H6}
\]
\end{lemma}
\begin{proof}
Write \(z=v-w\). Real linearity and (H2) give
\[
\begin{split}
\|z-\lambda Az\|^2
 &=\|z\|^2-2\lambda\langle Az,z\rangle+\lambda^2\|Az\|^2\\
 &\leq(1-2\lambda m+\lambda^2M^2)\|z\|^2\\
 &=(1-2\delta)\|z\|^2
 \leq(1-\delta)^2\|z\|^2 .
\end{split}
\]
Both comparison radii are nonnegative, so this proves (H6).
The finite rational proof expands the same squared inner product. Auxiliary
precision parameters account for evaluations of \(A\), and their errors are
removed by Lemma~\ref{lem:separation}. No spectral theorem, eigenvector
selection, orthonormal basis, or minimization principle is required.
\end{proof}

\subsection{A solution producer with all resource bounds}

Choose the explicit natural numbers
\[
\begin{split}
h&=\ceil{1/\delta},\qquad
R=\max(1,\ceil{\lambda B/\delta}),\qquad
C=\max(1,\ceil{2\lambda/\delta}),\\
k_n&=n+3+C,\qquad N_n=h(n+R+3),\qquad
\eta_n=\delta\,2^{-n-3}.
\end{split}
\tag{H7}
\]
For requested output precision \(n\), run the finite recurrence
\[
\begin{split}
 z^{\,n}_0&=0,\\
 z^{\,n}_{j+1}
   &=z^{\,n}_j-\lambda a(z^{\,n}_j,k_n)+\lambda b(k_n)
       \quad(j<N_n),\\
 c_n&=z^{\,n}_{N_n}.
\end{split}
\tag{H8}
\]
Every operation in (H7)--(H8) is primitive recursive on finite codes.
On malformed inputs or invalid positive constants the producer can
return \(0\); correctness is asserted on the valid domain.

\begin{theorem}[Quantitative operator Lax--Milgram]
\label{thm:hilbert-solver}
The sequence \(c_n\) produced by (H8) is regular and defines a unique
solution \(u\) of \(Au=b_\infty\). More precisely,
\[
\begin{gathered}
       \|c_n-c_\ell\|\leq2^{-n-1}\quad(\ell\geq n),\\
       \|c_n-u\|\leq2^{-n-2},\qquad
       \|u\|\leq B/m.
\end{gathered}
\tag{H9}
\]
The construction is uniform in the numerical inputs and schematic in the
fixed primitive-recursive operator and vector producers.
\end{theorem}
\begin{proof}
Let \(F_n(v)\) denote the finite update in (H8). Since \(2^C\geq C\),
\[
 \|F_n(v)-T(v)\|\leq2\lambda\,2^{-k_n}
                  \leq\delta2^{-n-3}=\eta_n.
 \tag{H10}
\]
Thus for any precision levels and inputs,
\[
\|F_n(v)-F_\ell(w)\|
 \leq q\|v-w\|+\eta_n+\eta_\ell .
\tag{H11}
\]
This is the central finite comparison contract.

First bound every intermediate state. Since
\(\|T(0)\|\leq\lambda B\leq\delta R\), induction gives
\[
\|z^{\,\ell}_j\|\leq R+\eta_\ell/\delta.
\tag{H12}
\]
Multiplication by \(q\), followed by addition of
\(\delta R+\eta_\ell\), preserves this bound.

Now fix \(\ell\geq n\), so \(N_\ell\geq N_n\).
Compare the \(n\)-run with the \emph{last \(N_n\) steps} of the \(\ell\)-run:
\[
             z^{\,n}_i
       \quad\hbox{and}\quad z^{\,\ell}_{N_\ell-N_n+i},
                    \qquad 0\leq i\leq N_n.
\]
Their initial distance is at most \(R+\eta_\ell/\delta\) by (H12).
Lemma~\ref{lem:perturb} and (H11) bound the final distance by
\[
\begin{split}
q^{N_n}(R+\eta_\ell/\delta)+(\eta_n+\eta_\ell)/\delta
 &\leq Rq^{N_n}+\eta_n/\delta+2\eta_\ell/\delta\\
 &\leq2^{-n-3}+2^{-n-3}+2\cdot2^{-\ell-3}\\
 &\leq2^{-n-1}.
\end{split}
\tag{H13}
\]
This proves regularity directly from two finite executions. The existence of
\(u\) was not assumed to prove convergence. The completion interpretation
of this regular name gives \(u\).
Taking arbitrarily fine comparison indices \(\ell\) in (H13) gives the sharper
\(\|c_n-u\|\leq Rq^{N_n}+\eta_n/\delta\leq2^{-n-2}\).
In the arithmetic reading this is a family of rational inequalities and
rational separation, not an added real limit rule.

To prove the equation using only comparisons of finite states, extend a run
by one more step and set \(w_j=\|z^{\,n}_{j+1}-z^{\,n}_j\|\).
By (H11),
\[
w_0\leq\lambda B+\eta_n,\qquad
w_{j+1}\leq q w_j+2\eta_n.
\]
Consequently
\[
\begin{split}
\|T(c_n)-c_n\|
 &\leq w_{N_n}+\eta_n\\
 &\leq q^{N_n}(\lambda B+\eta_n)+2\eta_n/\delta+\eta_n\\
 &\leq(2+3\delta)\,2^{-n-3}.
\end{split}
\tag{H14}
\]
The last bound is deliberately loose and tends to zero with an explicit
dyadic rate. Continuity of \(T\), the approximation estimate, and (H14)
show \(T(u)=u\). Since \(\lambda>0\), this is \(Au=b_\infty\).

If \(Au=Av\), coercivity applied to \(u-v\) gives
\(m\|u-v\|^2\leq0\), hence equality of the represented points.
Finally
\[
m\|u\|^2\leq\langle Au,u\rangle
 =\langle b_\infty,u\rangle\leq B\|u\|
\]
implies \(\|u\|\leq B/m\). This implication can be checked by rational upper
and lower tests; it does not require deciding whether the real norm is zero.
\end{proof}

\paragraph{Why this is an arithmetic proof plan.}
The state bound (H12), shifted two-run bound (H13), and adjacent-state bound
in (H14) are finite loop invariants. After replacing norms by rational upper
radii they are quantifier-free inequalities in state codes, loop counters,
constants, and chosen precision indices. A final observation point or test
vector is fixed before induction. The exact \(T\) in the notation is
eliminated using finite evaluators at an auxiliary precision; all additional
errors have explicit finite geometric-sum bounds. Separation yields (H11).
There is no induction on the formula ``\(z_j\) denotes a solution'' or on a
formula quantifying over the Hilbert space.
Section~\ref{sec:formalization} specifies the object-code obligations.

\begin{corollary}[Residual certification and stability]
\label{cor:residual}
For every represented \(v\),
\[
                   \|v-u\|\leq \frac{\|Av-b_\infty\|}{m}.
\tag{H15}
\]
Suppose a second operator \(\widetilde A\) has the same coercivity lower
bound \(m\), solution \(\widetilde u\), and data bound
\(\|\widetilde b\|\leq\widetilde B\). If
\[
\|b_\infty-\widetilde b\|\leq\beta,\qquad
\|(A-\widetilde A)x\|\leq\xi\|x\|,
\]
then
\[
       \|u-\widetilde u\|
       \leq \frac{\beta}{m}+\frac{\xi\widetilde B}{m^2}.
\tag{H16}
\]
\end{corollary}
\begin{proof}
For \(w=v-u\), coercivity and Cauchy--Schwarz give
\(m\|w\|^2\leq\langle Av-b_\infty,w\rangle
\leq\|Av-b_\infty\|\|w\|\).
The same nonnegative inequality argument as above proves (H15).
Apply it to \(\widetilde u\), whose residual under \(A\) is
\((A-\widetilde A)\widetilde u+\widetilde b-b_\infty\), and use
\(\|\widetilde u\|\leq\widetilde B/m\).
\end{proof}

Solving is therefore uniformly continuous on a stratum with a supplied
positive coercivity bound. The constants explain precisely why no uniform
extension to \(m=0\) is asserted.

\subsection{A real boundary-value problem, including variable coefficients}
\label{subsec:pde}

The following application constructs the strong operator instead of
postulating it. Let \(P=\Q[x]\) and
\(P_0=\{p\in\Q[x]:p(0)=p(1)=0\}\).
Equip \(P\) with inner product \(\int_0^1pq\), and \(P_0\) with
\[
                 \langle p,q\rangle_E=\int_0^1p'(x)q'(x)\,dx .
\]
All integrals are rational, since
\(\int_0^1x^j\,dx=1/(j+1)\).
These forms are positive definite. For a nonzero rational polynomial of degree
\(d\), one of \(d+1\) distinct rational interior points has nonzero value:
otherwise its coefficient vector is annihilated by an invertible Vandermonde
matrix. A coefficient bound gives an explicit Lipschitz constant, and hence
a rational interval on which its square has a positive rational lower bound.
Its squared integral is therefore positive. For the energy form, a polynomial
with zero derivative is constant, and the endpoint condition forces that
constant to be zero. This also gives finite arithmetic proofs of the
pre-Hilbert premises.
Let \(L^2_P\) and \(H_E\) denote the named completions.
These are the chosen presentations of the \(L^2\) and zero-boundary energy
spaces; no general measurable-set universe is needed.

\begin{lemma}[Polynomial inequalities and Green operator]
\label{lem:green}
For \(v\in P_0\),
\[
          \sup_{0\leq x\leq1}|v(x)|\leq\|v\|_E,\qquad
          \|v\|_{L^2}\leq\|v\|_E .
\tag{H17}
\]
For \(p\in P\) define
\[
(Gp)(x)=x\int_0^1(1-t)p(t)\,dt-\int_0^x(x-t)p(t)\,dt .
\tag{H18}
\]
Then \(Gp\in P_0\), \(-(Gp)''=p\), and
\[
  \langle Gp,v\rangle_E=\int_0^1p v,\qquad
  \|Gp\|_E\leq\|p\|_{L^2}.
\tag{H19}
\]
\end{lemma}
\begin{proof}
The finite polynomial identity \(v(x)=\int_0^x v'(t)\,dt\) and
Cauchy--Schwarz give
\(|v(x)|^2\leq x\int_0^x(v')^2\leq\|v\|_E^2\).
Integrating yields (H17). The integrals in (H18) are explicit polynomial
antiderivatives, so the endpoint and second-derivative assertions follow
by coefficient calculation. Polynomial integration by parts, with the
zero boundary values, proves the first equation in (H19). Taking \(v=Gp\),
\[
\|Gp\|_E^2=\int p\,Gp
 \leq\|p\|_{L^2}\|Gp\|_{L^2}
 \leq\|p\|_{L^2}\|Gp\|_E,
\]
which proves the norm estimate.

Cauchy--Schwarz for polynomial integrals follows by integrating the square
of \(sp+tq\) and using nonnegativity for every rational \(s,t\); the resulting
quadratic has the required discriminant inequality. Nonnegativity of the
integral of a polynomial known nonnegative on an interval follows from
finite Riemann sums and its explicit polynomial Lipschitz bound. Thus every
analytic-looking step has a finite rational-sum proof with a modulus, as in
the existing integration packet \cite{repoFramework}. The integration and
derivative identities themselves are exact finite coefficient identities.
\end{proof}

An energy name therefore defines a uniform continuous-function name:
(H17) bounds its uniform error by its energy error. The derivative need only
be an \(L^2\) name; a pointwise derivative is not inferred.

\begin{theorem}[A weak elliptic Dirichlet equation]
\label{thm:pde}
Suppose \(a\) is given by primitive-recursive rational polynomials \(a_k\)
with uniform error at most \(2^{-k}\) on \([0,1]\), and rational bounds
\[
                     0<m\leq a(x)\leq M .
\]
Suppose \(f\in L^2_P\) is given by a primitive-recursive polynomial name
\(f_k\), with \(\|f_k-f\|_{L^2}\leq2^{-k}\) and \(\|f\|_{L^2}\leq B\).
Then a primitive-recursive regular sequence of rational zero-boundary
polynomials names a unique \(u\in H_E\) such that
\[
                \int_0^1 a\,u'v'=\int_0^1f v
                         \qquad(v\in H_E).
\tag{H20}
\]
Its energy error at index \(n\) is at most \(2^{-n-2}\); the same bound
holds in uniform norm. Equation (H20) is the weak meaning of
\(-\left(a u'\right)'=f\), \(u(0)=u(1)=0\).
\end{theorem}
\begin{proof}
For \(v\in P_0\), define \(Av\) by its derivative
\[
 (Av)'=a v'-\int_0^1a(t)v'(t)\,dt,\qquad (Av)(0)=0.
\tag{H21}
\]
The derivative has mean zero, so also \((Av)(1)=0\).
Its energy interpretation is obtained by replacing \(a\) with a rational
polynomial approximation, integrating, and subtracting the mean.
The map \(g\mapsto g-\int_0^1g\) is contractive in \(L^2\), because
\[
       \int_0^1(g-\overline g)^2
                      =\int_0^1g^2-\overline g^{\,2}\leq\int_0^1g^2 .
\tag{H22}
\]
An approximation \(a_k\) consequently changes (H21) by at most
\(2^{-k}\|v\|_E\).
The rational \(U(v)=1+\|v\|_E^2\) bounds \(\|v\|_E\).
For accuracy \(2^{-r}\), take
\(k=r+\max(1,\lceil U(v)\rceil)\).
Then \(2^{-k}U(v)\leq2^{-r}\).
This constructs a primitive-recursive strong evaluator of \(A\).

For \(v,w\in P_0\), the mean in (H21) contributes nothing to
\(\int(Av)'w'\), because \(\int_0^1w'=0\). Hence
\[
 \langle Av,w\rangle_E=\int_0^1a v'w',\qquad
 \langle Av,v\rangle_E\geq m\|v\|_E^2,\qquad
 \|Av\|_E\leq M\|v\|_E.
\tag{H23}
\]
The last inequality follows from (H22) and
\(\|a v'\|_{L^2}\leq M\|v'\|_{L^2}\).
The explicit polynomial errors prove the same inequalities in the
completion. All strong-operator data have now been constructed.

Set \(b_\infty=Gf\) using (H18)'s extension. Equation (H19) gives the
strong evaluator \(Gf_k\), the required error, and \(\|b_\infty\|_E\leq B\).
Theorem~\ref{thm:hilbert-solver} produces \(Au=Gf\).
Pairing with \(v\) and using (H19), (H23) proves (H20), first for polynomials
and then for regular energy names by the displayed boundedness inequalities.
The same approximation argument defines the products and integrals in
(H20), with errors controlled by Cauchy--Schwarz.
Uniqueness and the uniform error follow from coercivity and (H17).
\end{proof}

\begin{example}[An exactly checkable continuous example]
\label{ex:pde}
Take \(a(x)=1+x\), \(f(x)=1\), \(m=1\), \(M=2\), \(B=1\).
Then \(\lambda=1/4\), \(\delta=1/8\), \(h=8\), \(R=2\), \(C=4\), and
\(N_n=8(n+5)\). Both \(A\) and \(Gf=x(1-x)/2\) have exact rational-polynomial
evaluators. Starting at \(0\), (H8) computes a polynomial of degree at most
\(N_n+1\). The classical expression for the same solution is
\[
                       u(x)=\frac{\log(1+x)}{\log 2}-x .
\]
Indeed \((1+x)u'(x)=1/\log2-(1+x)\), its derivative is \(-1\), and both
endpoints vanish. This expression is explanatory only: the polynomial
producer and its error proof do not depend on evaluating a logarithm.
The companion program checks rational energy residuals of the produced
polynomials and uses (H15) to certify them.
\end{example}

\subsection{Why weak functional data are insufficient}
\label{subsec:riesz}

The word \emph{strong} in Definition~\ref{def:strong-operator} matters.
A bilinear form given only by values on finite rational vectors need not
give a strong operator evaluator; similarly a bounded functional need not
give a norm-convergent name for its representing vector.

\begin{theorem}[Exact missing datum for a rational \(\ell^2\) vector]
\label{thm:tail-equivalence}
Let \(a_i\in\Q\) be a fixed primitive-recursive sequence whose finite squared
sums have a supplied common upper bound. The following data are equivalent:
\begin{enumerate}
\item a primitive-recursive strong \(\ell^2\) name with coordinates \(a_i\);
\item a primitive-recursive tail modulus \(\tau\) such that, for all \(k\)
and finite \(L\geq\tau(k)\),
\[
                    \sum_{\tau(k)\leq i<L}a_i^2\leq4^{-k}.
\tag{H24}
\]
\end{enumerate}
Both directions are explicit transformations without unbounded search.
\end{theorem}
\begin{proof}
Replace \(\tau(k)\) by \(\max_{j\leq k}\tau(j)\) if necessary.
The vector of the first \(\tau(k)\) coordinates is a regular name:
the squared distance between its \(k\)-th and \(\ell\)-th truncations,
\(\ell\geq k\), is bounded by (H24).
There is no need to assume \(\tau(k)\to\infty\). If a fixed coordinate is
still omitted, (H24), applied to the finite tail ending at that coordinate,
bounds its absolute value by \(2^{-k}\). Thus the represented coordinate
is \(a_i\) in either case.

Conversely, let \(p_k\) be a finite-support rational vector at distance
at most \(2^{-k}\) from the named vector. Choose \(\tau(k)\) larger than
every index in its support. On every finite coordinate interval starting at
\(\tau(k)\), \(p_k\) vanishes. Orthogonal projection onto that interval is
contractive, so the squared sum of the corresponding \(a_i\)'s is bounded
by \(4^{-k}\). The projection inequality says merely that a partial sum of
nonnegative rational squares is at most the full sum. For named coordinates
use finer approximants and rational separation. The support bound is
computed from the code \(p_k\), and hence is primitive recursive.
\end{proof}

\begin{proposition}[A bounded functional without a computable strong representative]
\label{prop:riesz-obstruction}
There is a rational primitive-recursive functional on finite rational
sequences, with norm bound \(1\), whose continuous extension to
\(\ell^2\) has no computable strong Riesz representative. In particular,
the missing tail modulus cannot be uniformly manufactured in \(\ipra\).
\end{proposition}
\begin{proof}
Fix a standard effective enumeration of machines and a primitive-recursive
pairing \(\langle e,t\rangle\). Whether machine \(e\) first halts at step
\(t\) is decided by a bounded simulation. Define
\[
 a_{\langle e,t\rangle}=
 \begin{cases}
 2^{-e-1},&\text{if \(e\) first halts at step \(t\)},\\
 0,&\text{otherwise.}
 \end{cases}
\tag{H25}
\]
This is a primitive-recursive rational sequence. Each \(e\) contributes
at most one nonzero square, so every finite squared sum is at most
\(\sum_{e\geq0}4^{-e-1}=1/3<1\). Only the finite geometric-sum inequality
is needed to prove this bound.

For finite rational \(v\), put \(L(v)=\sum_i a_i v_i\).
This is a primitive-recursive exact finite sum, and Cauchy--Schwarz gives
\(|L(v)|\leq\|v\|\). It extends to a computable functional on strong
\(\ell^2\) names: apply it to an input approximant and retain the same error.
A Riesz representative must have coordinate \(L(e_i)=a_i\).

A computable strong name for that representative would, by the proof of
Theorem~\ref{thm:tail-equivalence}, give a computable tail modulus.
For a machine \(e\), put \(N=\tau(e+2)\). A nonzero coordinate
\(\langle e,t\rangle\geq N\) would contribute
\(4^{-e-1}>4^{-e-2}\) to a finite tail, contradicting (H24).
Thus \(e\) halts exactly when one of the finitely many coordinates \(i<N\)
encoding \(e\) is nonzero. This decides the halting problem, a contradiction
by its elementary diagonal undecidability proof.
\end{proof}

This is a metatheoretic noncomputability result, not a formal independence
proof of a particular sentence of \(\ipra\).
It already blocks an unrestricted computable solution operator: take the
bilinear form \(\langle u,v\rangle\) and the functional \(L\). The weak
Lax--Milgram conclusion would require exactly the missing strong
representative. The positive PDE theorem avoids this problem by constructing
its Green representative and strong operator.

\begin{corollary}[Coercive bilinear values need not supply a strong operator]
\label{cor:bilinear-obstruction}
Even an exactly primitive-recursive bilinear form on finite rational vectors,
with coercivity bound \(1\) and boundedness constant \(2\), need not have
a computable strong representing operator.
\end{corollary}
\begin{proof}
Use the functional \(L\) just constructed and set
\(\mathfrak a(v,w)=\langle v,w\rangle+L(v)L(w)\).
This form is computed by rational finite sums. Its diagonal is at least
\(\|v\|^2\), and Cauchy--Schwarz and \(\|L\|\leq1\) give
\(|\mathfrak a(v,w)|\leq2\|v\|\|w\|\).
In the ordinary completion its representing operator is \(A=I+a\otimes a\),
where \(a\) is the vector of coordinates (H25).
Choose the index \(j\) of a known immediately halting machine and its known
halting time. Then \(a_j>0\) is a specified rational. A strong evaluator of
\(A\), applied to the finite basis vector \(e_j\), would give the computable
strong name
\[
                              a=(Ae_j-e_j)/a_j.
\]
This contradicts Proposition~\ref{prop:riesz-obstruction}.
\end{proof}
